\documentclass[11pt]{article}
\usepackage[T1]{fontenc}
\usepackage{lmodern}
\usepackage[margin=1in,headheight=15pt]{geometry}
\usepackage{amsmath,amssymb,amsthm,mathtools,mathrsfs}
\usepackage{booktabs,array,longtable,microtype,xurl,fancyhdr}
\usepackage[dvipsnames]{xcolor}
\usepackage[colorlinks=true,linkcolor=MidnightBlue,citecolor=MidnightBlue,urlcolor=MidnightBlue]{hyperref}
\usepackage{bookmark}
\emergencystretch=3em
\setlength{\parskip}{3pt}
\numberwithin{equation}{section}
\newtheorem{theorem}{Theorem}[section]
\newtheorem{lemma}[theorem]{Lemma}
\newtheorem{proposition}[theorem]{Proposition}
\newtheorem{corollary}[theorem]{Corollary}
\newtheorem{conjecture}[theorem]{Conjecture}
\theoremstyle{definition}\newtheorem{definition}[theorem]{Definition}
\theoremstyle{remark}\newtheorem{remark}[theorem]{Remark}
\newcommand{\dd}{\,\mathrm d}
\newcommand{\R}{\mathbb R}
\newcommand{\C}{\mathbb C}
\newcommand{\E}{\mathbb E}
\newcommand{\ii}{\mathrm i}
\DeclareMathOperator{\Li}{Li}
\DeclareMathOperator{\Impart}{Im}
\DeclareMathOperator{\Var}{Var}
\DeclareMathOperator{\Cov}{Cov}
\newcommand{\cut}{\C\setminus[1,\infty)}
\newcommand{\basecommit}{\nolinkurl{28357e8ca63dd78327db91d9be239d75e4462879}}
\pagestyle{fancy}\fancyhf{}
\fancyhead[L]{\small The real-order threshold for harmonic polylogarithms}
\fancyhead[R]{\small ProveIt}
\fancyfoot[C]{\thepage}
\hypersetup{pdftitle={The Sharp Real-Order Threshold for Harmonic Polylogarithms},pdfauthor={ProveIt Contributors},pdfsubject={Signed Hausdorff moments, endpoint atoms, angular zeros, and exact rational certificates}}
\title{\textbf{The Sharp Real-Order Threshold\\for Harmonic Polylogarithms}\\[7pt]
\large Endpoint atoms, zero geometry, and certified fractional-order identities}
\author{ProveIt Contributors}
\date{9 October 2026 (America/Los\_Angeles)\\[3pt]
\small Research continuation at a pinned repository revision}
\begin{document}
\maketitle\thispagestyle{empty}
\begin{abstract}
For $a,b>0$, let $F_{a,b}(z)=\sum_{n\ge1}H_{n-1}^{(b)}z^n/n^a$.
The existing ProveIt signed-kernel theory treats real $a\ge1$.
We prove that the exact boundary for a finite signed Hausdorff-moment
representation is instead $a+b\ge1$. Above this boundary an explicit
regularized Hurwitz kernel has precisely one negative--positive sign
change. On the boundary the density is replaced by a positive atom at
one and a strictly negative absolutely continuous part. This distinction
is necessary, not a matter of convention.

The construction extends slit-plane nonvanishing and the unique simple
angular zero to the full closed region $a+b\ge1$. On the critical line
we prove strict increase of the normalized displacement
$\cos\theta_{a,b}(\rho)/\rho$, through a positive integral formula for its
derivative; a dual left-endpoint theorem gives strict decrease for
$\Li_a(z)-z$. A positive two-variable representation proves disk
nonvanishing and Gaussian negativity even below the signed-measure
threshold. We derive elementary critical-kernel series, complete
monotonicity identities for a normalized Hurwitz expression, and strict
Hankel positivity for critical harmonic defects.

The Euler evaluator extends with explicit parameter-dependent bounds.
A standard-library-only rational certificate proves
$\Im F_{1/10,9/10}(\ii)<-0.52$, disproving extension of the old uniform
constant one beyond its stated domain. Eight value enclosures and twelve
angular brackets use integer-root inequalities and exact remainder
bounds. The existing $S_6$ identity and the broader normalized-radius
conjecture remain unproved here. The supplied symbolic checks and
numerical diagnostics are distinguished from the written proofs and
rational certificates throughout.
\end{abstract}
\clearpage
\begingroup\small
\tableofcontents
\endgroup
\clearpage

\section{Scope, baseline, and conventions}
\label{sec:scope}

\subsection{What is already known in the repository}
The baseline is the repository \emph{ProveIt}, revision \basecommit,
and in particular the canonical manuscript
\emph{Polylogarithms and their Arithmetic Bridges} and the continuation
\emph{Gaussian Polylogarithms Beyond Parity} \cite{repo,angular}.
The canonical reading artifact is described by its README as a 298-page,
twelve-chapter manuscript. Its current worklog marks $S_4$ as proved and
$S_6$ as conjectural. Those status distinctions are retained.

The signed-kernel chapter gives integer-order moment and zero theorems.
The angular continuation already extends the kernel and global
nonvanishing to real $a\ge1,b>0$, proves several parameter comparisons
at integer outer order, and proves the optimal universal Euler constant
one on the domain $a\ge1$. Re-proving those statements alone would not
constitute a continuation. Our main new domain is $0<a<1$, especially
the transition $a+b=1$. The endpoint measure, its consequences, and
the critical normalized-radius theorem are not present in the inspected
baseline. The classical Hurwitz regularization and hyperbolic partial
fractions used in the derivations are not claimed as new
\cite{dlmf-hurwitz,dlmf-hyperbolic}.

\begin{center}
\small
\begin{tabular}{@{}>{\raggedright\arraybackslash}p{.24\textwidth}>{\raggedright\arraybackslash}p{.31\textwidth}>{\raggedright\arraybackslash}p{.36\textwidth}@{}}
\toprule
Question & Inspected baseline & This article\\
\midrule
Signed moments & Real $a\ge1$, $b>0$ & Exact existence threshold $a+b\ge1$\\[3pt]
Critical endpoint & Not included in that domain & Required atom, exact mass, weak emergence\\[3pt]
Slit-plane zeros & None except zero for $a\ge1$ & Same conclusion for $a+b\ge1$\\[3pt]
Normalized radius & Broad conjecture remains & Critical-line and dual endpoint theorems\\[3pt]
Euler error & Optimal constant one for $a\ge1$ & Explicit extended bounds; certified scope obstruction\\[3pt]
Below threshold & Outside signed-kernel theorem & Positive double integral, disk nonvanishing\\
\bottomrule
\end{tabular}
\end{center}

The new results are proved below as ordinary mathematical theorems.
They have not been proof-assistant formalized or independently refereed.
A targeted source audit is not an exhaustive priority search: no claim
is made that every identity below is absent from the entire literature.
For general positive Hausdorff sequences and their analytic generating
functions, see Liu and Pego \cite{liu-pego}. Here the principal measure
is signed, and its one-crossing structure must be established separately.

\subsection{Definitions and boundary values}
For $a,b>0$, write
\begin{equation}\label{eq:def}
 c_n=c_n(a,b)=\frac{H_{n-1}^{(b)}}{n^a},\qquad
 H_m^{(b)}=\sum_{j=1}^m j^{-b},\qquad
 F_{a,b}(z)=\sum_{n=1}^{\infty}c_nz^n\quad(|z|<1).
\end{equation}
In particular, $c_1=0$ and $c_2=2^{-a}$. For integer indices,
$F_{a,b}(z)=\Li_{a,b}(z,1)$ in the decreasing-index convention.
Set $w=a+b$ and $g_{a,b}=\Impart F_{a,b}(\ii)$, where continuation
is on $\cut$. Section~\ref{sec:double} constructs that continuation
for every positive pair of parameters.

A signed Hausdorff representation means a finite real signed Borel
measure $\mu$ on $[0,1]$ such that
\begin{equation}\label{eq:moments}
 c_n=\int_{[0,1]}u^{n-1}\dd\mu(u),\qquad n\ge1.
\end{equation}
Thus $\mu([0,1])=0$. Equality of all these moments determines a finite
signed measure uniquely: polynomials are uniformly dense in
$C([0,1])$, and integration against a finite measure is continuous
in the uniform norm. The total variation is denoted by
$\|\mu\|_{\mathrm{TV}}$. Its two Jordan masses are equal, and we write
$m(\mu)=\mu_+([0,1])=\mu_-([0,1])$.

\emph{Boundary-series caution.} For $w>1$ the usual boundary series
converges at every $|z|=1,z\ne1$. On $w=1$ its coefficients tend to a
nonzero constant, and for $w<1$ they grow without bound. Accordingly,
in those latter cases $g_{a,b}$ and boundary identities mean analytic
or Abel values, not ordinary sums of a convergent series. This is proved
in Section~\ref{sec:threshold} and is essential for the Euler discussion.

\section{The regularized Hurwitz kernel}
\label{sec:regularization}

For $s>0$ and $T>0$ put $k_s(T)=T^{s-1}/\Gamma(s)$, and define
\begin{equation}\label{eq:h}
 h(t)=\frac1{1-e^{-t}}-\frac1t\qquad(t>0).
\end{equation}
This notation $h$ is local to this article; it is not the positive
primitive denoted by $h_{a,b}$ in the earlier angular report.

\begin{lemma}[Elementary regularizer]\label{lem:h}
The function $h$ increases strictly from $1/2$ to $1$, and
\begin{equation}\label{eq:hderivative}
 h'(t)=\frac1{t^2}-\frac1{4\sinh^2(t/2)}>0.
\end{equation}
For $b>0$, $b\ne1$, and $q>0$,
\begin{equation}\label{eq:hurwitz}
 \zeta(b,q)=\frac{q^{1-b}}{b-1}
 +\frac1{\Gamma(b)}\int_0^\infty e^{-qt}t^{b-1}h(t)\dd t.
\end{equation}
At $b=1$ the limiting formula is
\begin{equation}\label{eq:digamma}
 \int_0^\infty e^{-qt}h(t)\dd t=\log q-\psi(q).
\end{equation}
\end{lemma}
\begin{proof}
Since $2\sinh(t/2)>t$ for $t>0$, differentiation proves
\eqref{eq:hderivative}. Taylor expansion at zero and the definition
at infinity give the stated limits. For $b>1$, split the standard
Hurwitz integral kernel as
$(1-e^{-t})^{-1}=t^{-1}+h(t)$. The first integral is
$\Gamma(b-1)q^{1-b}$, proving \eqref{eq:hurwitz} there.
The remaining integral converges locally uniformly for complex
$\Re b>0$, because $h$ is bounded. Analytic continuation in $b$
proves the formula throughout that half-plane away from $b=1$.
This is the standard regularized Hurwitz formula, equivalently
DLMF 25.11.27 \cite{dlmf-hurwitz}. Taking the constant terms in
$\zeta(1+\varepsilon,q)=\varepsilon^{-1}-\psi(q)+O(\varepsilon)$
and $q^{-\varepsilon}/\varepsilon=\varepsilon^{-1}-\log q+O(\varepsilon)$
gives \eqref{eq:digamma}.
\end{proof}

Define the positive convolution
\begin{align}
 C_{a,b}(T)
 &=\frac1{\Gamma(a)\Gamma(b)}
   \int_0^T(T-t)^{a-1}t^{b-1}h(t)\dd t\label{eq:C}\\
 &=k_w(T)\E[h(TV)],\qquad V\sim\mathrm{Beta}(b,a).
 \label{eq:Cbeta}
\end{align}
The second parameter of this beta distribution is the outer order
$a$. This order of parameters matters. Immediate bounds are
\begin{equation}\label{eq:Cbounds}
 \tfrac12 k_w(T)<C_{a,b}(T)<k_w(T).
\end{equation}
Tonelli's theorem and the gamma integral give
\begin{equation}\label{eq:CLaplace}
 \int_0^\infty e^{-qT}C_{a,b}(T)\dd T
 =q^{-a}\left(\zeta(b,q)-\frac{q^{1-b}}{b-1}\right)
 \quad(b\ne1).
\end{equation}
For $b=1$ the right side is $q^{-a}(\log q-\psi(q))$.
All quantities in these formulas are finite for $q>0$.

\section{The sharp signed-measure threshold}
\label{sec:threshold}

\begin{theorem}[Exact existence threshold]\label{thm:threshold}
Let $a,b>0$. The sequence \eqref{eq:def} has a finite signed
Hausdorff representation if and only if $a+b\ge1$.

If $w>1$, the representing measure is absolutely continuous:
\begin{equation}\label{eq:mu-density}
 \dd\mu_{a,b}(u)=K_{a,b}(-\log u)\dd u\qquad(0<u<1),
\end{equation}
with no endpoint atoms. For $b\ne1$ its density in logarithmic
coordinates is
\begin{equation}\label{eq:K}
 K_{a,b}(T)=\zeta(b)k_a(T)-\frac{k_{w-1}(T)}{b-1}-C_{a,b}(T).
\end{equation}
For $b=1$ it is
\begin{equation}\label{eq:K1}
 K_{a,1}(T)=k_a(T)\bigl(\gamma+\psi(a)-\log T\bigr)-C_{a,1}(T).
\end{equation}
There is a unique $r\in(0,1)$ such that the density is strictly
negative on $(0,r)$ and strictly positive on $(r,1)$.

If $w=1$, necessarily $0<a,b<1$ and $a=1-b$. The unique measure is
\begin{equation}\label{eq:criticalmeasure}
 \mu_{a,b}=\frac1a\delta_1-\nu_{a,b},\qquad
 \dd\nu_{a,b}(u)=
 \bigl[-\zeta(b)k_a(-\log u)+C_{a,b}(-\log u)\bigr]\dd u.
\end{equation}
Here $\nu_{a,b}$ has strictly positive density, total mass $1/a$,
and no endpoint atoms. In particular $\|\mu_{a,b}\|_{\mathrm{TV}}=2/a$.
\end{theorem}

\subsection{Existence and the endpoint atom}
\begin{proof}[Proof of the representation and its necessity]
When $w>1$, every term of \eqref{eq:K} is integrable with weight
$e^{-T}\dd T$; this follows from $a>0$, $w-1>0$ and
\eqref{eq:Cbounds}. The logarithm in \eqref{eq:K1} is also integrable.
By the change of variables $u=e^{-T}$,
\[
 \int_0^1u^{n-1}K(-\log u)\dd u
 =\int_0^\infty e^{-nT}K(T)\dd T.
\]
Notice the factor $e^{-T}$ in the measure: $K(T)\dd T$ itself is
\emph{not} its logarithmic-coordinate measure.

Using \eqref{eq:CLaplace}, the Laplace transform of \eqref{eq:K} is
\[
 \zeta(b)n^{-a}-\frac{n^{-(w-1)}}{b-1}
 -n^{-a}\left(\zeta(b,n)-\frac{n^{1-b}}{b-1}\right)
 =\frac{\zeta(b)-\zeta(b,n)}{n^a}=c_n.
\]
The finite Hurwitz difference is exactly $H_{n-1}^{(b)}$.
At $b=1$, differentiating the gamma integral gives
\[
 \int_0^\infty e^{-nT}k_a(T)\log T\dd T
 =n^{-a}\bigl(\psi(a)-\log n\bigr).
\]
Together with \eqref{eq:digamma}, this makes the Laplace transform
of \eqref{eq:K1} equal to $n^{-a}(\gamma+\psi(n))=c_n$.
In particular the total mass, the $n=1$ case, is zero.

On the critical line, $\zeta(b)<0$ for $0<b<1$.
For example, the alternating Dirichlet series $\eta(b)>0$ and
$\zeta(b)=\eta(b)/(1-2^{1-b})$ prove this sign. Thus the density
in \eqref{eq:criticalmeasure} is strictly positive and integrable.
Because $w=1$, formula \eqref{eq:CLaplace} becomes
\[
 \int_0^\infty e^{-nT}C_{a,b}(T)\dd T
 =n^{-a}\zeta(b,n)+\frac1a.
\]
The moments of \eqref{eq:criticalmeasure} are consequently
\[
 \frac1a+\zeta(b)n^{-a}
 -\left(n^{-a}\zeta(b,n)+\frac1a\right)=c_n.
\]
The $n=1$ identity gives $\nu([0,1])=1/a$.
The positive atom and the negative density are mutually singular,
so the asserted total variation is exact.

It remains to exclude $w<1$. Then $0<b<1$ and the integral test gives
$H_{n-1}^{(b)}\sim n^{1-b}/(1-b)$, whence
$c_n\sim n^{1-w}/(1-b)\to\infty$. Moments of a finite signed measure
on $[0,1]$ have absolute value at most its total variation. This is
impossible. Uniqueness in the other cases follows from polynomial
density, as explained in Section~\ref{sec:scope}.
\end{proof}

\subsection{The one-crossing argument}
\begin{proof}[Proof of the strict sign pattern when $w>1$]
Divide the proposed density by the strictly positive $k_a(T)$ and
write $B(T)=K_{a,b}(T)/k_a(T)$. For $b\ne1$, differentiation of
\eqref{eq:Cbeta} yields
\begin{align}\label{eq:Bprime}
 B'(T)={}&-\frac{\Gamma(a)}{\Gamma(w-1)}T^{b-2}\notag\\
 &-\frac{\Gamma(a)}{\Gamma(w)}T^{b-1}
 \E\bigl[b h(TV)+TVh'(TV)\bigr]<0.
\end{align}
Both terms have the displayed sign. In particular the apparent
sign change of $1/(b-1)$ in the original density disappears after
normalization and differentiation. Differentiation under the beta
integral is valid on compact $T$-intervals since $h$ and $h'$ extend
continuously to zero.

For $b>1$, $B(T)\to\zeta(b)>0$ as $T\downarrow0$.
For $0<b<1$ and $w>1$, the term
$\Gamma(a)T^{b-1}/((1-b)\Gamma(w-1))$ makes that limit $+\infty$.
At infinity the term $-\Gamma(a)T^b\E h(TV)/\Gamma(w)$ dominates,
so $B(T)\to-\infty$. This proves a unique strict crossing.

For $b=1$,
\[
 B(T)=\gamma+\psi(a)-\log T-\frac{T}{a}\E h(TV),
 \qquad V\sim\mathrm{Beta}(1,a),
\]
and hence
\begin{equation}\label{eq:Bprime1}
 B'(T)=-\frac1T-\frac1a\E\bigl[h(TV)+TVh'(TV)\bigr]<0.
\end{equation}
Again the limits are $+\infty$ and $-\infty$.
In the $u$ coordinate, increasing $T$ means decreasing $u$;
therefore the crossing has precisely the negative--positive order
stated in the theorem. At the crossing $K'=k_aB'<0$, so the crossing
is transversal as well as unique.
\end{proof}

\begin{corollary}[Why the critical atom cannot be omitted]\label{cor:atom}
On $w=1$, every finite signed representing measure must satisfy
$\mu(\{1\})=1/a$. No integrable density alone can represent the sequence.
\end{corollary}
\begin{proof}
The same integral-test estimate gives $c_n\to1/a$ on the critical
line. Dominated convergence applied to $u^{n-1}$ against the total
variation of a finite measure gives
$\lim c_n=\mu(\{1\})$. An absolutely continuous measure has no such atom.
Equivalently, the limiting term $k_0$ in a formal convolution should
be a Dirac mass at $T=0$, not the zero function on $T>0$.
\end{proof}

\begin{corollary}[Ordinary versus Abel boundary values]\label{cor:boundary}
For $w>1$, the series defining $F_{a,b}$ converges ordinarily at every
$|z|=1,z\ne1$ and equals its integral continuation. If $w\le1$, that
ordinary boundary series does not converge at any such $z$.
\end{corollary}
\begin{proof}
For $w>1$, use the measure without endpoint atoms and the finite
geometric sum:
\[
 \sum_{n=1}^Nc_nz^n
 =\int\frac{z(1-(zu)^N)}{1-zu}\dd\mu(u).
\]
For a fixed boundary point other than one the denominator is bounded
away from zero uniformly in $0\le u\le1$. Dominated convergence
therefore applies. If $w=1$ the terms have modulus tending to $1/a$;
if $w<1$ their modulus grows without bound. In either case the
necessary term test fails.
\end{proof}

\section{Critical defects and the emergence of the atom}
\label{sec:defects}

\begin{theorem}[Completely monotone critical defects]\label{thm:defect}
Let $a=1-b\in(0,1)$ and
\[
 D_n=\frac1a-\frac{H_{n-1}^{(b)}}{n^a},\qquad n\ge1.
\]
Then $D_n=\int_0^1u^{n-1}\dd\nu_{a,b}(u)$ with the strictly positive
density in \eqref{eq:criticalmeasure}. If $\Delta D_n=D_n-D_{n+1}$,
then $\Delta^kD_n>0$ for every $k\ge0,n\ge1$. Moreover,
\begin{equation}\label{eq:defectbounds}
 -\zeta(b)n^{-a}+\frac1{2n}<D_n
 <-\zeta(b)n^{-a}+\frac1n.
\end{equation}
Every finite Hankel matrix $(D_{j+k+1})_{j,k=0}^r$ is positive
definite, and $D_nD_{n+2}>D_{n+1}^2$.
\end{theorem}
\begin{proof}
Subtract the moments of \eqref{eq:criticalmeasure} from its atomic
mass. Then
\[
 \Delta^kD_n=\int_0^1u^{n-1}(1-u)^k\dd\nu(u)>0.
\]
On $w=1$, $k_w(T)=1$, so \eqref{eq:Cbounds} reads
$1/2<C_{a,b}(T)<1$. Integrating these inequalities against
$e^{-nT}\dd T$ gives \eqref{eq:defectbounds}.
For a nonzero real vector $(v_0,\ldots,v_r)$,
\[
 \sum_{j,k=0}^rv_jv_kD_{j+k+1}
 =\int_0^1\left(\sum_{j=0}^rv_ju^j\right)^2\dd\nu(u)>0,
\]
since a nonzero polynomial cannot vanish on an interval and the
density is positive everywhere inside $(0,1)$. Strict log-convexity
follows either from its $2\times2$ minors or from strict
Cauchy--Schwarz with the weight $u^{n-1}\dd\nu$.
\end{proof}

The theorem also gives a continuous version. For $q>0$,
\begin{equation}\label{eq:continuousdefect}
 D(q)=\frac1a-q^{-a}\bigl(\zeta(b)-\zeta(b,q)\bigr)
 =\int_0^\infty e^{-qT}
   \bigl[-\zeta(b)k_a(T)+C_{a,b}(T)\bigr]\dd T.
\end{equation}
Thus $(-1)^jD^{(j)}(q)>0$ for all integers $j\ge0$.
All derivatives are justified by exponential domination. These are
inequalities for specific harmonic/Hurwitz expressions, not claims
about independence of their special values.

\begin{theorem}[Weak emergence of the critical atom]\label{thm:weak}
Fix $0<b<1$ and let $a_\varepsilon=1-b+\varepsilon$ with
$\varepsilon\downarrow0$. Then
$\mu_{a_\varepsilon,b}\Rightarrow\mu_{1-b,b}$ weakly on $[0,1]$.
The corresponding $F$ converge locally uniformly on $\cut$.
\end{theorem}
\begin{proof}
Put $M=1/(1-b)$. Formula \eqref{eq:K} can be written
\[
 \dd\mu_{a_\varepsilon,b}(u)
 =M k_\varepsilon(-\log u)\dd u-\dd\nu_\varepsilon(u),
\]
where the density of $\nu_\varepsilon$ is
$-\zeta(b)k_{a_\varepsilon}(-\log u)+C_{a_\varepsilon,b}(-\log u)>0$.
The first term has logarithmic-coordinate law $M$ times a
$\mathrm{Gamma}(\varepsilon,1)$ density. Its mean $T$ is
$\varepsilon$, so it converges weakly to $M\delta_1$ in the
$u=e^{-T}$ coordinate.

The negative densities converge pointwise on $(0,1)$ to the critical
negative density. Their total masses are all $M$, because the positive
term has mass $M$ and the signed measure has total mass zero. The
limiting negative measure also has mass $M$. Scheff\'e's lemma
therefore gives convergence of the negative densities in $L^1$.
For completeness, this lemma follows from
$\int|f_n-f|=\int f_n+\int f-2\int\min(f_n,f)$ and dominated
convergence for $\min(f_n,f)\le f$.

This proves weak convergence with uniformly bounded total variation.
The kernels $z/(1-zu)$ are continuous in $u$ and uniformly bounded
with their $z$-derivatives on compact subsets of the slit plane.
A finite-net argument then upgrades pointwise convergence of their
integrals to locally uniform convergence.
\end{proof}

\section{A positive two-variable representation at all positive orders}
\label{sec:double}

The threshold in Theorem~\ref{thm:threshold} concerns a
\emph{one-variable finite signed measure}. It is not a threshold for
analytic continuation or every positivity property.

\begin{theorem}[Positive double kernel]\label{thm:double}
For every $a,b>0$ and $z\in\cut$,
\begin{equation}\label{eq:double}
 F_{a,b}(z)=z^2\int_{0<y<x<1}
 \frac{W_{a,b}(x,y)}{(1-xz)(1-yz)}\dd y\dd x,
\end{equation}
where
\begin{equation}\label{eq:W}
 W_{a,b}(x,y)=\frac{(-\log x)^{a-1}(\log(x/y))^{b-1}}
 {\Gamma(a)\Gamma(b)}.
\end{equation}
The weight is strictly positive and has total mass $2^{-a}$.
\end{theorem}
\begin{proof}
For $|z|<1$, expand the two denominators. Substitution $y=xv$
gives, for nonnegative integers $j,k$,
\[
 \int_{0<y<x<1}W_{a,b}(x,y)x^jy^k\dd y\dd x
 =\frac1{(j+k+2)^a(k+1)^b}.
\]
Summing over $j+k=n-2$ yields $c_n$. Absolute convergence justifies
the rearrangement, and taking $j=k=0$ gives the total mass.
On each compact subset of $\cut$, both denominators are bounded
away from zero uniformly in $x,y\in[0,1]$. The integral is consequently
holomorphic there and supplies the stated continuation.
\end{proof}

\begin{corollary}[Disk nonvanishing below and above the threshold]
\label{cor:disk}
For all $a,b>0$, $F_{a,b}$ has no nonzero zeros in
$\{z:|z|\le1,\ z\ne1\}$. More precisely, when $\Impart z>0$ and
$|z|\le1$,
\begin{equation}\label{eq:diskpick}
 \Impart\frac{F_{a,b}(z)}z>0.
\end{equation}
At any upper-half-disk zero of $\Impart F_{a,b}(z)$, the value
$F_{a,b}(z)$ itself is strictly negative real.
\end{corollary}
\begin{proof}
A direct calculation gives
\[
 \Impart\frac{z}{(1-xz)(1-yz)}
 =\frac{(\Impart z)(1-xy|z|^2)}{|1-xz|^2|1-yz|^2}>0
\]
for $0<y<x<1$, proving \eqref{eq:diskpick}. The lower half follows
by conjugation. On the real interval $z<1$, the integral
$F(z)/z^2$ is strictly positive. Finally, if $F(z)=R\in\R$ with
$\Impart z>0$, then $\Impart(R/z)=-R\Impart z/|z|^2>0$, so $R<0$.
\end{proof}

\begin{corollary}[Gaussian negativity and strict log-convexity]
\label{cor:gaussian}
For every $a,b>0$,
\begin{equation}\label{eq:gaussianpositive}
 -g_{a,b}=\int_{0<y<x<1}W_{a,b}(x,y)
       \frac{x+y}{(1+x^2)(1+y^2)}\dd y\dd x>0,
\end{equation}
and
\begin{equation}\label{eq:gaussianbounds}
 \frac{H_2^{(b)}}{4\,3^a}<-g_{a,b}<\frac{H_2^{(b)}}{3^a}.
\end{equation}
The function $(a,b)\mapsto\Gamma(a)\Gamma(b)(-g_{a,b})$ is strictly
log-convex on the positive quadrant.
\end{corollary}
\begin{proof}
Set $z=\ii$ in \eqref{eq:double} and rationalize the denominator.
Since $1<(1+x^2)(1+y^2)<4$ and
$\int W(x+y)=c_3=H_2^{(b)}/3^a$, the bounds follow.
For log-convexity apply H\"older's inequality to the positive integral
with the gamma factors removed. Its parameter dependence is the
exponential of
$(a-1)\log(-\log x)+(b-1)\log\log(x/y)$.
Equality in H\"older would require a nontrivial real linear combination
of these two logarithms to be constant almost everywhere on the open
triangle, which is impossible. This proves strictness for every
nontrivial line segment of parameter pairs.
\end{proof}

\subsection{All-order integral identities}
\begin{theorem}[Hurwitz--polylogarithm transport]\label{thm:transport}
For $a,b>0$, $b\ne1$, and $z\in\cut$,
\begin{align}\label{eq:transport}
 F_{a,b}(z)={}&\zeta(b)\Li_a(z)-\frac{\Li_{a+b-1}(z)}{b-1}\notag\\
 &-\frac1{\Gamma(b)}\int_0^\infty t^{b-1}h(t)\Li_a(ze^{-t})\dd t.
\end{align}
At $b=1$ the pole-free identity is
\begin{equation}\label{eq:transport1}
 F_{a,1}(z)=\gamma\Li_a(z)-\partial_a\Li_a(z)
       -\int_0^\infty h(t)\Li_a(ze^{-t})\dd t.
\end{equation}
In particular, the index $a+b-1$ in \eqref{eq:transport} is allowed
to be zero or negative.
\end{theorem}
\begin{proof}
Inside the disk, multiply \eqref{eq:hurwitz}, with $q=n$, by
$n^{-a}z^n$ and sum. Exponential decay at infinity, boundedness of
$h$, and $b>0$ justify termwise integration. The Hurwitz difference
is $H_{n-1}^{(b)}$, so coefficient comparison gives
\eqref{eq:transport}. At $b=1$, use
$H_{n-1}=\gamma+\log n-\int_0^\infty e^{-nt}h(t)\dd t$;
the sum of $(\log n)n^{-a}z^n$ is $-\partial_a\Li_a(z)$.
The integral expressions on the right are locally holomorphic on the
slit plane. Near $t=0$, $\Li_a(ze^{-t})$ stays bounded on compact
slit sets, and at infinity it is $O(e^{-t})$. Analytic continuation
therefore proves the identities throughout the claimed domain.
Standard conventions for $\Li_s$ are those of DLMF 25.12
\cite{dlmf-polylog}.
\end{proof}

These identities are useful even when no signed Hausdorff measure
exists. They are transformations to a depth-one polylogarithm integral,
not assertions that a depth-two constant has been reduced to a finite
basis of depth-one periods.

\section{Slit-plane nonvanishing and the unique angular zero}
\label{sec:zeros}

\begin{lemma}[Ordered-sign test]\label{lem:sign}
Suppose a nonzero finite signed measure $\mu$ has total mass zero,
its negative part is supported to the left of $r$, and its positive
part to the right, allowing an atom at $r$ in either part. If $q$ is
strictly increasing on the relevant supports and integrable against
$|\mu|$, then $\int q\dd\mu>0$. For strictly decreasing $q$ the
inequality is reversed. These conclusions apply both to the
one-crossing density of Theorem~\ref{thm:threshold} and to its critical
endpoint measure.
\end{lemma}
\begin{proof}
Subtract $q(r)\mu([0,1])=0$. On each nonzero signed part,
$(q(u)-q(r))\dd\mu(u)$ has the asserted sign. The strict positive
density on an interior interval ensures that the integral is strict,
even when the other part is concentrated at the endpoint $r=1$.
\end{proof}

\begin{theorem}[Zero-free principal branch on the full closed domain]
\label{thm:zerofree}
For $a,b>0$ with $w\ge1$, $F_{a,b}$ has exactly one zero on
$\cut$: its double zero at the origin. More explicitly, let $r$ be
the crossing of Theorem~\ref{thm:threshold} for $w>1$, and let $r=1$
for $w=1$. Then
\begin{align}
 F_{a,b}(z)&=\frac{z^2}{1-rz}P_{a,b}(z),\label{eq:factor}\\
 P_{a,b}(z)&=\int_{[0,1]}\frac{u-r}{1-zu}\dd\mu_{a,b}(u).
 \label{eq:P}
\end{align}
The measure $(u-r)\dd\mu_{a,b}(u)$ is positive and nonzero.
The imaginary part of $P_{a,b}(z)$ has the strict sign of
$\Impart z$, and $P_{a,b}(x)>0$ for real $x<1$.
\end{theorem}
\begin{proof}
The moment representation gives
$F(z)=\int z(1-zu)^{-1}\dd\mu(u)$, first in the disk and then on
the slit plane. Subtract $(1-rz)^{-1}\int z\dd\mu=0$ to obtain
\eqref{eq:factor}. Positivity of $(u-r)\dd\mu$ follows from the sign
pattern; on the critical line the atom disappears after multiplication
by $u-1$, leaving $(1-u)\dd\nu(u)>0$.
For nonreal $z$,
\[
 \Impart P(z)=(\Impart z)
   \int\frac{u(u-r)}{|1-zu|^2}\dd\mu(u),
\]
and the integral is strictly positive. For real $z<1$ all denominators
in \eqref{eq:P} are positive. Thus $P$ never vanishes there, and
$1-rz$ has no zero off the excluded cut. Finally
$P(0)=\int u\dd\mu(u)=c_2=2^{-a}$, proving exact multiplicity two.
\end{proof}

The factorization mechanism is already used in the angular continuation
for $a\ge1$ \cite{angular}. The substantive extension here is the
construction of the required signed measure for $0<a<1,w\ge1$,
including its atomic endpoint.

\begin{theorem}[Unique simple angular zero]\label{thm:angular}
Let $a,b>0$, $w\ge1$, and $0<\rho\le1$. There is exactly one
$\theta_{a,b}(\rho)\in(0,\pi)$ satisfying
$\Impart F_{a,b}(\rho e^{\ii\theta})=0$. It obeys
\begin{equation}\label{eq:anglerange}
 \arccos\rho<\theta_{a,b}(\rho)<\frac\pi2.
\end{equation}
The imaginary part is positive before this zero and negative after
it, and its angular derivative at the zero is strictly negative.
The zero depends smoothly on the radius, including a one-sided
interpretation at $\rho=1$.
\end{theorem}
\begin{proof}
Set $c=\cos\theta$ and
\[
 K_{\rho,c}(u)=\frac1{1-2\rho cu+\rho^2u^2},\qquad
 J(\rho,c)=\int K_{\rho,c}(u)\dd\mu(u).
\]
Then $\Impart F(\rho e^{\ii\theta})=\rho\sin\theta J(\rho,c)$.
The weight $K_{\rho,0}$ is strictly decreasing, so $J(\rho,0)<0$.
For $\rho<1$, the weight $K_{\rho,\rho}$ is strictly increasing on
$(0,1)$, since its derivative has the sign of $1-u$. Thus
$J(\rho,\rho)>0$, providing a zero $0<c<\rho$.

For $\rho=1$ and $w>1$, the negative measure is supported away from
$u=1$. Separating the two signs before expanding the endpoint kernel
gives
\[
 \lim_{c\uparrow1}J(1,c)=\int\frac{\dd\mu(u)}{(1-u)^2}
 =\sum_{n\ge1}n c_n>0,
\]
with $+\infty$ allowed. The negative part is integrable at this
endpoint and monotone convergence applies to the positive part.
On $w=1$ that particular argument must not be used: the negative
density also reaches one. Instead write $\mu=M\delta_1-\nu$.
Then
\[
 (1-c)J(1,c)=\frac M2
 -\int\frac{1-c}{1-2cu+u^2}\dd\nu(u)\longrightarrow\frac M2>0.
\]
Indeed the last integrand tends to zero at every $u<1$ and is
bounded by one for $0<c<1$, because its denominator is at least
$1-c^2\ge1-c$. The finite measure $\nu$ has no atom at one.
This proves existence at unit radius as well.

For $c'>c$,
\begin{equation}\label{eq:kernelratio}
 \frac{\dd}{\dd u}\frac{K_{\rho,c'}(u)}{K_{\rho,c}(u)}
 =\frac{2\rho(c'-c)(1-\rho^2u^2)}
 {(1-2\rho c'u+\rho^2u^2)^2}>0\quad(0<u<1).
\end{equation}
At a zero $J(\rho,c)=0$, the measure $K_{\rho,c}\dd\mu$ has zero
mass and the same ordered signs. Lemma~\ref{lem:sign}, applied to
this ratio, gives $J(\rho,c')>0$ for every $c'>c$, and the reversed
inequality for every $c'<c$. This proves uniqueness and the global
sign pattern. At the zero,
\[
 J_c=\int K_{\rho,c}(u)\frac{2\rho u}{1-2\rho cu+\rho^2u^2}
       \dd\mu(u)>0,
\]
since the second factor is strictly increasing on $(0,1)$.
Consequently the angular derivative is
$-\rho\sin^2\theta J_c<0$. All differentiations are dominated away
from the point $z=1$, and the implicit function theorem gives smooth
radial dependence.
\end{proof}

By Theorem~\ref{thm:weak} and the strict sign brackets in this proof,
for every fixed $0<\rho\le1$ and $0<b<1$,
\[
 \theta_{1-b+\varepsilon,b}(\rho)
 \longrightarrow\theta_{1-b,b}(\rho).
\]
Thus the atom is compatible with continuous angular geometry; it
changes the representing measure and summation interpretation, not the
existence of a well-defined limiting root.

\section{Normalized-radius monotonicity at opposite endpoints}
\label{sec:radius}

The normalized radial quantity is
\begin{equation}\label{eq:eta}
 \eta_{a,b}(\rho)=\frac{\cos\theta_{a,b}(\rho)}{\rho}.
\end{equation}
The earlier report conjectures its decrease for integer $a\ge2$ and
all $b>0$, and for $a=1,b>1$, but its increase for $a=1,0<b<1$
\cite[Conjecture ``Normalized-radius monotonicity'']{angular}.
No such broad statement follows from uniqueness of the zero.
We prove two endpoint-measure theorems that isolate an exact
mechanism for the two directions.

\subsection{A right-endpoint atom}
\begin{theorem}[Critical normalized-radius theorem]\label{thm:criticalradius}
Let $a=1-b\in(0,1)$. The function $\eta_{a,b}(\rho)$ is strictly
increasing for $0<\rho\le1$, and
\begin{equation}\label{eq:criticalrange}
 \frac12<\eta_{a,b}(\rho)<1,
 \qquad
 \arccos\rho<\theta_{a,b}(\rho)<\arccos(\rho/2).
\end{equation}
The following positive integral gives its derivative. Put
$x=\rho^2$, $m=2\eta-1$, $\dd\omega(u)=(1-u)\dd\nu_{a,b}(u)$,
and
$D(u)=1+x\{u^2-(1+m)u\}$. Then $0<m<1$ and
\begin{equation}\label{eq:criticalderivative}
 \frac{\dd m}{\dd x}
 =\frac{\displaystyle\int_0^1
           \frac{(u-m)^2(1-u)}{D(u)^2}\dd\omega(u)}
 {(1-xm)\displaystyle\int_0^1
           \frac{1-xu}{D(u)^2}\dd\omega(u)}>0.
\end{equation}
In particular $\eta'(\rho)=\rho m'(\rho^2)>0$.
\end{theorem}
\begin{proof}
We prove the identity for any measure $M\delta_1-\nu$ with
$\nu([0,1])=M$ and strictly positive density on $(0,1)$, which includes
\eqref{eq:criticalmeasure}. With $c=\rho(1+m)/2$,
$D(1)=1-xm$ and
\[
 J(\rho,c)=\frac{M}{D(1)}-\int\frac{\dd\nu(u)}{D(u)}
 =\frac{x}{1-xm}G(m,x),
\]
where
\begin{equation}\label{eq:criticalG}
 G(m,x)=\int\frac{m-u}{D(u)}\dd\omega(u).
\end{equation}
Here we used
$D(u)-D(1)=x(u-1)(u-m)$.
For fixed $0<x\le1$, $G(0,x)<0$ and $G(1,x)>0$.
At $x=1,m=1$, the latter integrand simplifies to $\dd\nu(u)$,
so it is still integrable and positive. It is also the left limit
as $m\uparrow1$: for $1/2\le m<1$, the corresponding scalar
integrand against $\nu$ lies between $-1$ and $1$, so dominated
convergence applies. For interior $0<m<1$,
\begin{equation}\label{eq:Gm}
 G_m=\int\frac{1-xu}{D(u)^2}\dd\omega(u)>0.
\end{equation}
Thus its root lies strictly in $(0,1)$ and coincides with the
angular root already proved unique.

Let $d(u)=u^2-(1+m)u$. At a root of $G$, the probability measure
proportional to $\dd\omega(u)/D(u)$ has mean $m$. The elementary
identities
\[
 d(m)=-m,\qquad d(u)-d(m)=(u-m)(u-1)
\]
therefore give
\begin{align*}
 G_x
 &=\int (u-m)\frac{d(u)}{D(u)}\frac{\dd\omega(u)}{D(u)}\\
 &=\int (u-m)\left(\frac{d(u)}{D(u)}
                         -\frac{-m}{1-xm}\right)
                         \frac{\dd\omega(u)}{D(u)}\\
 &=-\frac1{1-xm}\int\frac{(u-m)^2(1-u)}{D(u)^2}\dd\omega(u)<0.
\end{align*}
The middle subtraction integrates to zero. Combining
$m'=-G_x/G_m$ with \eqref{eq:Gm} proves
\eqref{eq:criticalderivative}. All denominators are bounded away from
zero locally at the root, including at $x=1$ because $m<1$.
Strictness follows from the nondegenerate positive density.
\end{proof}

\begin{corollary}[A positive small-radius coefficient]\label{cor:smallradius}
Let $a=1-b\in(0,1)$ and abbreviate
$A=H_2^{(b)}$, $B=H_3^{(b)}$, $C=H_4^{(b)}$. Then
\begin{align}\label{eq:smallradius}
 \eta_{a,b}(\rho)={}&\frac A2\left(\frac23\right)^a\notag\\
 &+\left\{AB\left(\frac13\right)^a
   -\frac C2\left(\frac25\right)^a
   -\frac{A^3}{2}\left(\frac8{27}\right)^a\right\}\rho^2
 +O(\rho^4).
\end{align}
The expression in braces is strictly positive on this critical line.
More intrinsically, put
$m_0=(c_3-c_2)/c_2$ and let $U$ have probability measure
$\dd\omega/c_2$. That coefficient equals
\begin{equation}\label{eq:variance}
 \frac12\E[(U-m_0)^2(1-U)]>0.
\end{equation}
\end{corollary}
\begin{proof}
At $x=0$, $G(m,0)=0$ gives $m_0=\int u\dd\omega/\int\dd\omega$.
The moments of $\omega$ are
\[
 \int u^j\dd\omega=D_{j+1}-D_{j+2}=c_{j+2}-c_{j+1}.
\]
Thus its mass is $c_2$ and
$(1+m_0)/2=c_3/(2c_2)=A(2/3)^a/2$.
Formula \eqref{eq:criticalderivative} at $x=0$ gives the coefficient
\eqref{eq:variance}. Expressing its moments through $c_2,c_3,c_4,c_5$
and simplifying gives
\[
 \frac{c_3c_4}{c_2^2}-\frac{c_5}{2c_2}
       -\frac{c_3^3}{2c_2^3},
\]
which is exactly the braces in \eqref{eq:smallradius}.
Analytic implicit dependence near $x=0$ gives the remainder.
The formal expansion itself is already present for $a\ge1$ in the
angular report; the positive variance formula and the strict sign
on $a+b=1$ are the additional conclusions here.
\end{proof}

\subsection{A dual left-endpoint atom}
\begin{theorem}[Opposite monotonicity for a left atom]\label{thm:left}
Suppose $\mu=\nu-M\delta_0$, where $\nu$ has strictly positive
integrable density on $(0,1)$ and mass $M$. Define
$F(z)=\int z(1-zu)^{-1}\dd\mu(u)$.
On every upper semicircle of radius $0<\rho\le1$, its imaginary
part has one simple zero, and its normalized displacement
$\eta(\rho)=\cos\theta(\rho)/\rho$ satisfies
$0<\eta<1/2$ and $\eta'(\rho)<0$.

More explicitly, set $x=\rho^2$, $m=2\eta$, $\dd\omega=u\dd\nu$,
and $D(u)=1+x(u^2-mu)$. At the root,
\begin{equation}\label{eq:leftderivative}
 \frac{\dd m}{\dd x}
 =-\frac{\displaystyle\int\frac{u(u-m)^2}{D(u)^2}\dd\omega(u)}
 {\displaystyle\int\frac1{D(u)^2}\dd\omega(u)}<0.
\end{equation}
\end{theorem}
\begin{proof}
In this case
\[
 J(\rho,c)=\int\frac{\dd\nu(u)}{D(u)}-M=xG(m,x),
 \quad
 G(m,x)=\int\frac{m-u}{D(u)}\dd\omega(u).
\]
The function $G$ is negative at $m=0$, positive at $m=1$, and
\[
 G_m=\int D(u)^{-2}\dd\omega(u)>0,
 \qquad
 G_x=\int u(u-m)^2D(u)^{-2}\dd\omega(u)>0.
\]
The denominators stay positive for $0\le m\le1$ and $0\le x\le1$.
This proves the unique simple root, its range, and
\eqref{eq:leftderivative}. The signs outside the root also follow
from the same increasing-kernel-ratio argument as before.
\end{proof}

\begin{corollary}[A polylogarithm-minus-linear-term family]\label{cor:Li}
For every real $a>0$, the normalized angular displacement of
$\Li_a(z)-z$ decreases strictly with $\rho\in(0,1]$ and belongs to
$(0,1/2)$. This family is the limit of $F_{a,b}$ as $b\to\infty$.
\end{corollary}
\begin{proof}
The probability measure $\dd\nu(u)=k_a(-\log u)\dd u$ has moments
$n^{-a}$. Thus $\nu-\delta_0$ has moments $0$ at $n=1$ and $n^{-a}$
for $n\ge2$, which gives $\Li_a(z)-z$. Apply
Theorem~\ref{thm:left}. Coefficient convergence gives the indicated
limit in the disk. For sufficiently large $b>1$, the signed measures
of Theorem~\ref{thm:threshold} have uniformly bounded variation,
as proved below. Moment uniqueness and polynomial approximation then
give weak convergence to $\nu-\delta_0$, hence locally uniform
convergence on the slit plane.
\end{proof}

These two theorems do not settle the finite-$b$ normalized-radius
conjecture from the baseline. They show why the orientation of an
endpoint atom can force opposite directions, and supply rigorous
boundary data for a continuous-parameter phase diagram.

\section{Elementary critical-kernel identities and Hurwitz inequalities}
\label{sec:criticalidentities}

In this section $0<b<1$, $a=1-b$, and
$C_b(T)=C_{1-b,b}(T)=\E[h(TV)]$ with
$V\sim\mathrm{Beta}(b,1-b)$. In particular $1/2<C_b(T)<1$ and
$C_b$ is strictly increasing.

\begin{theorem}[An elementary convergent series]\label{thm:elementary}
For every $T>0$,
\begin{equation}\label{eq:elementary}
 C_b(T)=\frac12+\sum_{k=1}^\infty
 \frac{\sin\!\bigl(b\arctan(T/(2\pi k))\bigr)}
 {\pi k\bigl(1+(T/(2\pi k))^2\bigr)^{b/2}}.
\end{equation}
Every summand is positive. The tail after $K\ge1$ terms is bounded by
\begin{equation}\label{eq:elementarytail}
 0<C_b(T)-\frac12-\sum_{k=1}^K(\text{summand})
 \le\frac{bT}{2\pi^2K}.
\end{equation}
For $|T|<2\pi$ it also has the convergent expansion
\begin{equation}\label{eq:Bernoulli}
 C_b(T)=\frac12+\sum_{j=1}^\infty
 \frac{B_{2j}(b)_{2j-1}}{(2j)!(2j-1)!}T^{2j-1},
\end{equation}
where $(b)_m$ is the rising factorial and $B_{2j}$ is a Bernoulli
number, not a Bernoulli polynomial evaluated at $b$.
\end{theorem}
\begin{proof}
The classical partial fraction for $\coth$ gives
\begin{equation}\label{eq:hpartial}
 h(t)=\frac12+2t\sum_{k=1}^\infty\frac1{t^2+(2\pi k)^2}.
\end{equation}
This follows from DLMF 4.36.3 after rescaling
\cite{dlmf-hyperbolic}. The summands are positive for $t>0$, so
Tonelli permits averaging over $V$.
For $|z|<1$, the beta moments yield
\[
 \E\frac1{1-zV}=\sum_{j\ge0}z^j\frac{(b)_j}{j!}=(1-z)^{-b}.
\]
Both sides are analytic for $z\in\cut$, so the same identity holds
there. At $z=\ii q$, its imaginary part is
\[
 \E\frac{qV}{1+q^2V^2}
 =\frac{\sin(b\arctan q)}{(1+q^2)^{b/2}}.
\]
Taking $q=T/(2\pi k)$ proves \eqref{eq:elementary}.
Since $\sin(b\arctan q)\le bq$ and the denominator factor is at
least one, its tail is at most
$bT(2\pi^2)^{-1}\sum_{k>K}k^{-2}\le bT/(2\pi^2K)$.

The usual Bernoulli expansion of \eqref{eq:h} has radius $2\pi$:
$h(t)=1/2+\sum_{j\ge1}B_{2j}t^{2j-1}/(2j)!$.
For $|T|<2\pi$ it is uniformly absolutely convergent for
$0\le V\le1$. Average termwise and use
$\E V^{2j-1}=(b)_{2j-1}/(2j-1)!$ to obtain
\eqref{eq:Bernoulli}.
\end{proof}

For instance,
\begin{equation}\label{eq:Cexample}
 C_{1/2}(T)=\frac12+\frac{T}{24}-\frac{T^3}{2304}
                       +\frac{T^5}{122880}+O(T^7).
\end{equation}
The companion program regenerates exact coefficients through degree
$13$ at $b=1/10,1/2,9/10$ by two algebraically equivalent constructions.
The numerical tests of \eqref{eq:elementary} use independent beta
quadrature, with the analytic tail bound displayed separately.

\begin{theorem}[A completely monotone normalized Hurwitz expression]
\label{thm:hurwitzcm}
For $0<b<1$ and $q>0$, define
\begin{equation}\label{eq:L}
 L_b(q)=q^{b-1}\zeta(b,q)+\frac1{1-b}.
\end{equation}
Then
\begin{equation}\label{eq:criticalLaplace}
 L_b(q)=\int_0^\infty e^{-qT}C_b(T)\dd T,
 \qquad \frac1{2q}<L_b(q)<\frac1q.
\end{equation}
In particular, $(-1)^jL_b^{(j)}(q)>0$ for all $j\ge0$.
The function
\begin{equation}\label{eq:normalizedHurwitz}
 qL_b(q)=q^b\zeta(b,q)+\frac q{1-b}
\end{equation}
decreases strictly from $1$ to $1/2$ as $q$ increases from zero to
infinity.
\end{theorem}
\begin{proof}
Formula \eqref{eq:criticalLaplace} is \eqref{eq:CLaplace} with
$a=1-b$. Bounds and strict complete monotonicity follow from
$1/2<C_b<1$ and differentiating the positive Laplace integral.
Substitution $T=s/q$ gives
\[
 qL_b(q)=\int_0^\infty e^{-s}C_b(s/q)\dd s.
\]
Strict increase of $C_b$ proves strict decrease with $q$.
The endpoint limits follow by dominated convergence and the limits
$C_b(0+)=1/2$, $C_b(\infty)=1$.
\end{proof}

The beta average is the common source of the elementary series and
the Hurwitz inequalities. The inequalities, finite-difference signs,
and Hankel determinants are analytic consequences of a positive
integral, not numerical pattern extrapolations.

\section{Euler enclosures and the failure of an unrestricted constant one}
\label{sec:euler}

Let $a,b>0$ with $w\ge1$. Define
\begin{equation}\label{eq:eulerdef}
 f(n)=c_{2n+1}=\frac{H_{2n}^{(b)}}{(2n+1)^a},\qquad
 d_k=\sum_{n=0}^k(-1)^n\binom kn f(n),\qquad
 E_N=\sum_{k=0}^{N-1}\frac{d_k}{2^{k+1}}.
\end{equation}
The error is measured against $g_{a,b}=\Impart F_{a,b}(\ii)$.
On the critical line this is an Abel value; the ordinary alternating
series with terms $(-1)^nf(n)$ diverges by the term test.

\begin{theorem}[Extended one-sided Euler bounds]\label{thm:euler}
For every $k\ge1$ one has $d_k<0$, while $d_0=0$. The $E_N$ decrease
strictly to $g_{a,b}$, and
\begin{equation}\label{eq:exacteuler}
 E_N-g_{a,b}=-2^{-N}\int_{[0,1]}
       \frac{(1-u^2)^N}{1+u^2}\dd\mu_{a,b}(u)>0.
\end{equation}
Consequently
\begin{equation}\label{eq:eulerenclosure}
 E_N-M(a,b)2^{-N}<g_{a,b}<E_N\qquad(N\ge1),
\end{equation}
with the following elementary choices:
\begin{equation}\label{eq:Mchoices}
 M(a,b)=
 \begin{cases}
 1,&a\ge1,\\
 (1-b)^{-1},&0<a<1,\ 0<b<1,\ w\ge1,\\
 b/(b-1),&0<a<1,\ b>1,\\
 1+1/a,&0<a<1,\ b=1.
 \end{cases}
\end{equation}
The choice one in the first line is the prior optimal universal
constant on that restricted domain; it is not asserted to work in
the other lines.
\end{theorem}
\begin{proof}
The binomial theorem and \eqref{eq:moments} give
$d_k=\int(1-u^2)^k\dd\mu(u)$. Thus $d_0=0$, while the strictly
decreasing weight gives $d_k<0$ for $k\ge1$. With
$m=m(\mu)$, we have $|d_k|\le m$. Therefore the Euler series is
absolutely convergent and may be summed under the measure integral:
\[
 \sum_{k\ge0}\frac{d_k}{2^{k+1}}
 =\int\frac{\dd\mu(u)}{1+u^2}=\Impart F(\ii).
\]
Summing the geometric tail gives \eqref{eq:exacteuler}. Its weight
is strictly decreasing, bounded by one, and strictly less than one
for every $u>0$. Since the negative Jordan part has no atom at zero,
$0<E_N-g<m2^{-N}$. It remains to bound $m$.

For $b<1,w>1$, rewrite \eqref{eq:K} as
\[
 K_{a,b}=\frac{k_{w-1}}{1-b}
          -\bigl[-\zeta(b)k_a+C_{a,b}\bigr].
\]
The two displayed densities are positive, and each has mass
$1/(1-b)$ in the $u$ measure, by zero total mass. Hence
$m<1/(1-b)$. On $w=1$ the exact Jordan mass is $1/a=1/(1-b)$.
For $b>1$, \eqref{eq:K} is a difference of positive densities
$\zeta(b)k_a$ and $k_{w-1}/(b-1)+C_{a,b}$, so
$m<\zeta(b)<1+1/(b-1)=b/(b-1)$.

For $b=1$ and $0<a<1$, use a translation estimate rather than the
singular limit of either of these bounds. Let $\sigma_a$ be the
probability law of $u=e^{-Y}$ with $Y\sim\mathrm{Gamma}(a,1)$.
Its logarithmic density $p_a(s)=e^{-s}s^{a-1}/\Gamma(a)$ decreases
strictly. If $D_{e^{-t}}$ denotes pushforward under $u\mapsto e^{-t}u$,
then direct comparison of $p_a(s)$ and $p_a(s-t)$ gives
\[
 \|\sigma_a-D_{e^{-t}}\sigma_a\|_{\mathrm{TV}}
 =2\int_0^tp_a(s)\dd s
 \le\frac{2t^a}{\Gamma(a+1)}.
\]
The signed-measure integral
\[
 \widetilde\mu=\int_0^\infty
       \frac{\sigma_a-D_{e^{-t}}\sigma_a}{e^t-1}\dd t
\]
converges in total variation. Its $(n-1)$st moment is
$n^{-a}\int_0^\infty(1-e^{-(n-1)t})(e^t-1)^{-1}\dd t=c_n$.
By uniqueness it is $\mu_{a,1}$, and
\[
 m\le\frac1{\Gamma(a+1)}\int_0^\infty\frac{t^a}{e^t-1}\dd t
 =\zeta(a+1)<1+1/a.
\]

Finally, for $a\ge1$ the earlier sharper bound can be recovered
without any fractional-order gap. At $a=1$ the equivalent base density is
\[
 K_{1,b}(T)=-\frac{T^b}{\Gamma(b+1)}
 +\frac1{\Gamma(b)}\int_T^\infty\frac{t^{b-1}}{e^t-1}\dd t.
\]
Its $n$th Laplace value is
$-n^{-b-1}+n^{-1}\sum_{j=1}^n j^{-b}=H_{n-1}^{(b)}/n$,
by Tonelli and a finite geometric sum. Moment uniqueness therefore
identifies this base density with \eqref{eq:K} or \eqref{eq:K1}.
Both positive terms in this difference have $e^{-T}\dd T$ integral
one. For the second term, interchange the integrals and use
$(1-e^{-t})/(e^t-1)=e^{-t}$; the first is a gamma integral.
Their positive overlap makes the Jordan mass strictly less than one.
For $a>1$, convolution by $k_{a-1}$ in the $T$ coordinate preserves
moments after multiplication by $n^{-(a-1)}$ and contracts the
$L^1(e^{-T}\dd T)$ norm, because
$\int_0^\infty e^{-T}k_{a-1}(T)\dd T=1$.
Moment uniqueness identifies the result with the present kernel,
so $m<1$ for every real $a\ge1$. This proves all four choices.
\end{proof}

The same translation proof gives, when $0<a\le1$ and $a+b>1$,
the optional non-elementary bound
\[
 m\le\frac{\Gamma(a+b)\zeta(a+b)}{\Gamma(a+1)\Gamma(b)}.
\]
It may improve a particular estimate, but diverges as $a+b\downarrow1$.
The critical atomic construction is therefore not replaced by this
translation argument.

\begin{proposition}[Exact finite Euler weights]\label{prop:weights}
For $N\ge1$,
\begin{equation}\label{eq:weights}
 E_N=2^{-N}\sum_{n=0}^{N-1}(-1)^nf(n)
                  \sum_{j=n+1}^N\binom Nj.
\end{equation}
Thus the finite sum can be formed without constructing a triangular
difference table.
\end{proposition}
\begin{proof}
Substitute the definition of $d_k$ and interchange the finite sums.
The coefficient identity is
\[
 \sum_{k=n}^{N-1}2^{N-k-1}\binom kn=\sum_{j=n+1}^N\binom Nj,
\]
which follows by induction from Pascal's identity. Harmonic numbers
and binomial tails can then each be updated successively. This
rearrangement is already used in the integer-order manuscript; the
new feature in the implementation here is certified rational
interval arithmetic for fractional powers.
\end{proof}

\begin{theorem}[A certified scope obstruction]\label{thm:counterexample}
For $a=1/10$ and $b=9/10$,
\begin{align}\label{eq:counterinterval}
 -0.5297619322857712630946293655
 &<g_{1/10,9/10}\notag\\
 &<-0.5297619322857712630946293653.
\end{align}
In particular, $g_{1/10,9/10}<-13/25<-1/2$.
Since $E_1=0$, the inequality $E_N-2^{-N}<g_{a,b}$ fails at $N=1$
for this pair. A uniform constant one cannot be extended to the
entire region $a+b\ge1$. The failure also occurs strictly above
the threshold, with ordinary boundary convergence:
\begin{align}\label{eq:counterordinary}
 -0.5308442529661593261720712602
 &<g_{1/10,1}\notag\\
 &<-0.5308442529661593261720712600<-0.53.
\end{align}
\end{theorem}
\begin{proof}
The accompanying standard-library program \texttt{code/certify.py}
evaluates \eqref{eq:weights} with $N=96$ by outward rational interval
arithmetic. Every fractional power is bounded by integer tenth-root
inequalities, as detailed in Section~\ref{sec:certificates}. The
critical mass bound is the exact rational $M=10$. Subtracting
$10\,2^{-96}$ from the lower endpoint of the finite Euler enclosure
and leaving its upper endpoint unchanged gives the interval recorded
in \texttt{data/exact\_certificates.json}. Outward decimal rounding
yields \eqref{eq:counterinterval}; exact rational comparisons verify
its upper endpoint is less than $-13/25$. No numerical polylogarithm,
floating-point root, or unproved identity enters this certificate.
For $(a,b)=(1/10,1)$ the same computation uses $M=1+1/a=11$
and proves \eqref{eq:counterordinary}. Its ordinary boundary
convergence follows from $a+b=11/10>1$ and
Corollary~\ref{cor:boundary}.
\end{proof}

This is \emph{not} a counterexample to the baseline theorem, which
explicitly assumes $a\ge1$. It is a counterexample to dropping that
hypothesis while retaining the same universal error constant.

\begin{proposition}[The exact critical constant for fixed parameters]
\label{prop:fixedconstant}
On $a+b=1$, the scaled errors $2^N(E_N-g_{a,b})$ decrease strictly
to zero for $N\ge1$. The smallest constant in the non-strict estimate
$E_N-g_{a,b}\le C2^{-N}$, valid for all integers $N\ge1$ at that fixed
pair, is $C=-2g_{a,b}$.
\end{proposition}
\begin{proof}
The endpoint atom contributes zero to \eqref{eq:exacteuler} when
$N\ge1$. Hence
\[
 2^N(E_N-g)=\int_0^1\frac{(1-u^2)^N}{1+u^2}\dd\nu(u).
\]
Strict pointwise decrease and dominated convergence prove the first
claim. Its largest value occurs at $N=1$ and is $2(E_1-g)=-2g$.
For the corresponding strict inequality the infimum of admissible
constants is the same, but equality at $N=1$ prevents attaining it.
\end{proof}

\section{Exact certificates and independent diagnostics}
\label{sec:certificates}

\subsection{Certified fractional powers}
For a positive rational exponent $p/q$ and an integer $n\ge1$, set
$B=10^{65}$ and compute the integer
$r=\lfloor(n^pB^q)^{1/q}\rfloor$. The verifier explicitly checks
\begin{equation}\label{eq:integerroot}
 r^q\le n^pB^q<(r+1)^q.
\end{equation}
If equality holds on the left, $n^{-p/q}=B/r$ exactly. Otherwise
\begin{equation}\label{eq:rootinterval}
 \frac B{r+1}<n^{-p/q}<\frac B r.
\end{equation}
All operations in these statements are integer or rational. The
integer Newton iteration is merely a method to propose $r$;
\eqref{eq:integerroot}, not its convergence history, validates the
result. Additions, sign-aware scalar products, and positive products
are rounded outward onto the grid $10^{-60}\mathbb Z$ using integer
division. This keeps numerator sizes manageable without introducing
floating-point assumptions.

Finite Euler centers at fractional indices are generally algebraic,
not rational. The program encloses them by rational intervals, then
adds the proved rational tail. It does not mislabel an irrational
finite center as an exact rational value.

\begin{center}
\small
\begin{tabular}{@{}c c r@{}}
\toprule
$a$ & $b$ & $g_{a,b}$, displayed approximately\\
\midrule
$1/10$ & $9/10$ & $-0.52976193228577126309$\\
$1/10$ & $1$ & $-0.53084425296615932617$\\
$1/2$ & $1/2$ & $-0.40393698771701357983$\\
$3/4$ & $1/2$ & $-0.33913987183657185082$\\
$1/2$ & $1$ & $-0.40169999261081928166$\\
$1/4$ & $3/2$ & $-0.47502596533036419882$\\
$1$ & $1$ & $-0.27219826128795026631$\\
$2$ & $3$ & $-0.09324795037117472982$\\
\bottomrule
\end{tabular}
\end{center}
These displayed approximations are summaries, not the certified
endpoints. The JSON stores outward lower and upper decimal bounds
for each of the eight $N=96$ enclosures, with widths and mass bounds.
For $(1,1)$ the value agrees with the known
$-\pi\log2/8$; that identity is a convention check, not a premise of
the fractional-order counterexample.

\subsection{Angular certificates without trigonometric rounding}
Let $0<\rho<1$ and $-1<c<1$ be rational. The Chebyshev identity gives
\begin{equation}\label{eq:cheb}
 \frac{\Impart F_{a,b}(\rho e^{\ii\theta})}{\sin\theta}
 =\sum_{n\ge1}c_n\rho^n U_{n-1}(c),\qquad c=\cos\theta,
\end{equation}
where $U_{-1}=0$, $U_0=1$, and
$U_{n+1}=2cU_n-U_{n-1}$. Thus every polynomial value is rational.
For $|c|<1$, $|U_{n-1}(c)|\le n$; this follows from
$|\sin(n\theta)|\le n|\sin\theta|$. Also $c_n\le n$ for every
positive $a,b$. With $j=N+1$, a valid explicit tail is therefore
\begin{equation}\label{eq:angletail}
 \left|\sum_{n>N}c_n\rho^nU_{n-1}(c)\right|
 \le\rho^j\left\{\frac{j^2}{1-\rho}
       +\frac{2j\rho}{(1-\rho)^2}
       +\frac{\rho(1+\rho)}{(1-\rho)^3}\right\}.
\end{equation}
The formula sums $\sum_{n\ge j}n^2\rho^n$ exactly.

The numerical program proposes rational endpoints for $c$, but the
separate rational verifier accepts a bracket only when the interval
in \eqref{eq:cheb} is strictly negative at the lower endpoint and
strictly positive at the upper endpoint, including \eqref{eq:angletail}.
The theorem supplies uniqueness when $a+b\ge1$. For the deliberately
included test $(a,b,\rho)=(1/4,1/4,1/2)$, it certifies only existence
of a root in the bracket; no uniqueness conclusion is attached.

\begin{center}
\small
\begin{tabular}{@{}c c r r r@{}}
\toprule
$a$ & $b$ & $\eta(1/4)$ & $\eta(1/2)$ & $\eta(3/4)$\\
\midrule
$1/10$ & $9/10$ & $0.7387056020$ & $0.7429273646$ & $0.7517703302$\\
$1/2$ & $1/2$ & $0.6980111081$ & $0.7015741355$ & $0.7088461987$\\
$9/10$ & $1/10$ & $0.6719709290$ & $0.6751086399$ & $0.6814072072$\\
\bottomrule
\end{tabular}
\end{center}
The critical rows illustrate Theorem~\ref{thm:criticalradius}; the
proof does not infer monotonicity from this table. The underlying
rational $c$-brackets have width $5\times10^{-12}$. Nine critical
brackets, two strictly supercritical brackets, and one subcritical
existence bracket are replayed.

\subsection{Independent tests and their limits}
The recorded run has the following evidence layers.

\paragraph{Exact arithmetic.}
The standard-library verifier passes all eight Euler enclosures and
twelve angular brackets, checking 7,430 integer-root inequalities.
A symbolic program verifies nine rational identities underlying the
angular kernels, the two derivative formulas, and the small-radius
variance reduction. A separate symbolic routine checks 21 exact
Bernoulli coefficients at three rational parameters. None of these
finite checks certifies the analytic convergence arguments by itself.

\paragraph{Numerical diagnostics.}
Independent beta--gamma quadrature checks 50 Laplace moments; the
largest observed absolute discrepancy in the supplied run is about
$2.42\times10^{-13}$. Positive double quadrature separately evaluates
three critical Gaussian values. Twelve elementary-series tests compare
the partial sum and its analytic tail with beta quadrature. Twelve
critical Laplace-identity tests have largest observed discrepancy
about $2.54\times10^{-11}$. A 15-point critical-constant grid provides
conjectural evidence only. These are floating-point diagnostics, not
rigorous interval enclosures, and the stated discrepancies are
observations, not proved error bounds for the quadratures.

\paragraph{Reproduction.}
From the package root, with Python and the optional diagnostic
dependencies installed, run
\begin{verbatim}
python code/audit.py
python code/critical_identities.py
python code/certify.py
python code/build_article.py
\end{verbatim}
The third command alone replays the exact certificates using only
Python's standard library and the supplied rational root proposals.
The first two commands are not prerequisites for trusting the rational
endpoint tests, and floating-point root proposals are never accepted
without a fresh sign certificate. The build command requires a local
LaTeX installation. A separate manifest verifies package bytes; it
is an integrity check, not a mathematical proof.

\section{Audit conclusions and manuscript integration}
\label{sec:audit}

\subsection{Corrections that an extension actually requires}
No error in the \emph{stated} $a\ge1$ zero-free or Euler-constant-one
theorems was established during this targeted audit. The following
points are concrete corrections to an unrestricted continuation of
them, and should be made explicit when integrating the new results.

\paragraph{The critical kernel is a measure, not just a function.}
Substituting $a+b=1$ into a formula valid for $a+b>1$ and discarding
the $k_{a+b-1}$ term would delete the positive atom. The remaining
density has mass $-1/a$, not zero, and produces the wrong moments.
The correction is exactly $a^{-1}\delta_1$. Corollary~\ref{cor:atom}
also proves that no alternative integrable density can repair this.

\paragraph{Boundary convergence changes at the same line.}
At $a+b=1$, $c_n\to1/a$; the Gaussian alternating terms likewise
fail to tend to zero. Statements about ordinary convergence must
retain $a+b>1$. Euler acceleration on the critical line computes
the analytic/Abel value, with the measure proof justifying that value.
It must not be described as summing an ordinarily convergent series.

\paragraph{The uniform constant one does not survive automatically.}
Theorem~\ref{thm:counterexample} proves the obstruction with exact
rational arithmetic. The old hypothesis $a\ge1$ is mathematically
substantive. Use the bounds in \eqref{eq:Mchoices}, or a separately
proved sharper constant, in the extended region.

\paragraph{Fractional integration needs its own sign proof.}
A generic fractional integral does not come with an automatic
one-crossing theorem. The proof here uses the specific monotonicity
of $K/k_a$ in \eqref{eq:Bprime} and \eqref{eq:Bprime1}; it does not
invoke an unproved variation-diminishing assertion. Similarly,
slit-plane zero-freeness below $a+b=1$ must not be inferred from a
finite signed measure that does not exist.

\paragraph{Source status is not reset by a new delivery.}
The inspected repository has multiple historical $S_4$ certificates
and later integrated proofs. They remain proved prior material.
The $S_6$ identity remains conjectural here; we neither certify it
by numerical agreement nor assert its independence from a proposed
period basis. The broad normalized-radius conjecture remains open
outside the new regimes proved in Section~\ref{sec:radius}.

\subsection{Integration map}
The suggested destination for this complete package is
\nolinkurl{Analysis/Polylogarithms/docs/reports/real-order-threshold/}.
The main threshold and kernel formulas belong next to
\nolinkurl{chapters/05-signed-kernels.tex}; the critical endpoint and
its root geometry should precede any claim extending the angular
report beyond $a\ge1$. The elementary critical series and Hurwitz
inequalities can be cross-referenced from the integration and
differentiation chapters.

The supplied integration fragment uses the label prefix
\texttt{realorder:}. It is an editorial insertion with a summary of
the proved results, not a substitute for the complete proofs. The
package includes a claim-status ledger, a provenance file with pinned
source blob identifiers, and instructions for replaying checks. No
upstream source, historical certificate, or repository branch was
modified in preparing this delivery.

\section{Further research questions and precise conjectures}
\label{sec:future}

\subsection{The still-open normalized-radius problem}
The most direct task is to settle the existing conjecture for integer
$a\ge2$ and finite $b>0$, together with the $a=1$ sign transition at
$b=1$. Theorems~\ref{thm:criticalradius} and \ref{thm:left} supply
opposite endpoint mechanisms but do not interpolate between them.
A promising approach is to express the radial derivative at the root
as a covariance plus a controlled remainder. The right-atom proof
works because $d(u)-d(m)$ factors as $(u-m)(u-1)$; the left-atom proof
works because it factors as $u(u-m)$. For a general one-crossing
density no comparable fixed endpoint factor has yet been established.

For real outer orders a richer phase diagram is natural. The critical
line has increasing normalized radius, while the $b\to\infty$ family
has decreasing normalized radius for every $a>0$. A sharp description
should locate the transition in $(a,b)$ and decide whether the sign
can change with $\rho$ for one fixed pair. The small-radius coefficient
in \eqref{eq:smallradius}, valid as a formal local coefficient more
generally, gives a necessary local test, but does not prove a global
radial sign.

\subsection{Below-threshold zeros}
Theorem~\ref{thm:threshold} rules out finite signed Hausdorff measures
when $a+b<1$, not zero-freeness of the function. The positive double
kernel proves nonvanishing in the closed unit disk for all positive
orders. It also gives existence of an angular zero on each circle
of radius $0<\rho<1$: the imaginary part, divided by
$\rho^2\sin\theta$, has integrand
\[
 W_{a,b}(x,y)
 \frac{2c-\rho(x+y)}
 {(1-2\rho cx+\rho^2x^2)(1-2\rho cy+\rho^2y^2)}.
\]
At $c=0$ this is negative and at $c=\rho$ positive. It has fixed
sign outside $0<c<\rho$, so every such zero belongs to that interval.
Uniqueness is the additional issue.

\begin{conjecture}[Angular uniqueness below the moment threshold]
\label{conj:subcritical}
For $a,b>0$ with $a+b<1$, the imaginary part of $F_{a,b}$ has exactly
one simple zero on each upper semicircle of radius $0<\rho\le1$,
with the same sign pattern as in Theorem~\ref{thm:angular}.
\end{conjecture}
The supplied subcritical rational bracket proves existence at one
specific parameter triple only. It is not evidence for an exhaustive
zero count. At unit radius, even the positive endpoint argument must
be supplied separately; the interior-radius sign test above does not
justify an interchange at the singular endpoint $z=1$.

A second question is whether the slit-plane zero-free theorem holds
for all positive $a,b$. The factor $F/z$ in
Corollary~\ref{cor:disk} need not retain a positive imaginary part
outside the unit disk: the numerator $1-xy|z|^2$ can change sign.
A different positive representation or a counterexample is needed.
We leave this as a question rather than label it a proved extension.

\subsection{The optimal uniform critical Euler constant}
\begin{proposition}[Endpoint values for the critical error constant]
\label{prop:endlimits}
On $b=1-a$,
\begin{align}
 \lim_{a\downarrow0}(-2g_{a,1-a})&=\frac\pi4+\frac{\log2}{2},
 \label{eq:leftlimit}\\
 \lim_{a\uparrow1}(-2g_{a,1-a})&=\frac\pi2-1.
 \label{eq:rightlimit}
\end{align}
\end{proposition}
\begin{proof}
In \eqref{eq:gaussianpositive}, substitute
$x=e^{-S/2}$ and $y=xe^{-Y}$. The weight becomes $2^{-a}$ times
the product of independent $\mathrm{Gamma}(a,1)$ and
$\mathrm{Gamma}(b,1)$ probability densities for $S,Y$.
The remaining bounded continuous function is
$(x+y)/((1+x^2)(1+y^2))$.
As $a\downarrow0$, $S\to0$ in probability and $Y$ converges in law
to an exponential variable. Thus $x\to1$ and $y$ is asymptotically
uniform on $(0,1)$, giving
\[
 \lim(-g)=\int_0^1\frac{1+y}{2(1+y^2)}\dd y
 =\frac\pi8+\frac{\log2}{4}.
\]
As $a\uparrow1$, $Y\to0$, while $x$ has limiting density $2x$ on
$(0,1)$. Therefore
\[
 \lim(-g)=2\int_0^1\frac{x^2}{(1+x^2)^2}\dd x
 =\frac\pi4-\frac12.
\]
Weak convergence and boundedness justify both limits, proving the
formulas after multiplication by two.
\end{proof}

\begin{conjecture}[Sharp critical constant]\label{conj:constant}
The function $a\mapsto-2g_{a,1-a}$ decreases strictly on $(0,1)$.
Consequently the sharp constant uniform over the entire critical line
and all $N\ge1$ would be
\begin{equation}\label{eq:constantconjecture}
 C_{\mathrm{critical}}=\frac\pi4+\frac{\log2}{2}.
\end{equation}
\end{conjecture}
The 15-point diagnostic grid is consistent with this statement.
The endpoint limits and Proposition~\ref{prop:fixedconstant} prove
that any such uniform constant is at least the value in
\eqref{eq:constantconjecture}; the matching global upper bound is
not proved in this article. The stronger strict monotonicity is
also explicitly conjectural.

\subsection{More precise critical computation}
The elementary series \eqref{eq:elementary} has a simple positive
$O(K^{-1})$ tail. Its summands have expansions in inverse even powers
of $k$, so subtracting a finite number of those terms and restoring
zeta tails should produce arbitrarily accelerated, sign-controlled
versions. An exact implementation should retain outward bounds for
trigonometric and fractional-power factors rather than certify a
floating-point result after the fact. The Bernoulli expansion is
convergent only for $|T|<2\pi$ and should not be used as a global
large-$T$ expansion.

The sharp $N$-dependence of critical Euler errors is another concrete
target. Formula \eqref{eq:exacteuler} localizes at $u=0$, where the
negative density has both a gamma-logarithmic term and a beta-averaged
term. Uniform asymptotics as $a\downarrow0$ would have to resolve
competition between $N$ and $a$, not merely fix the parameters before
letting $N$ tend to infinity.

\subsection{Higher depth, affine steps, and formal proof}
A natural extension replaces $H_{n-1}^{(b)}$ by higher-depth nested
harmonic sums. The coefficient growth suggests candidate total-order
thresholds, but growth alone proves only nonexistence on one side;
it does not produce a one-crossing signed measure on the other.
Endpoint atoms may be replaced by several ordered signed components,
so angular uniqueness may need to be replaced by a finite zero-count
statement.

The manuscript's affine coefficients $H_{rn}^{(b)}/(dn+1)^a$ are
another useful direction. Their existing theorem covers $a\ge1$.
A sharp fractional-order threshold should depend on the growth
exponent, but the relative scales $r,d$ can affect the kernel sign
geometry and hence the best Euler constants. The cancellation between
a gamma endpoint and a Hurwitz tail needs to be redone rather than
assumed to survive unchanged.

Finally, the proof has a reasonably small formal core: gamma and beta
integrals, the regularized Hurwitz identity, polynomial density,
ordered-sign measures, and the two rational derivative factorizations.
A Lean formalization of the generic endpoint-atom root lemmas would
be reusable independently of polylogarithms. The current rational
verifier is a replayable finite computation, not such a formalization.

\section*{Conclusion}
\addcontentsline{toc}{section}{Conclusion}
The condition $a\ge1$ is sufficient but not sharp for the signed
harmonic-polylogarithm kernel. The exact finite-measure domain is
$a+b\ge1$, and its boundary carries a mathematically necessary atom.
That atom provides both a clean normalized-radius theorem and a warning:
ordinary boundary convergence and the old universal Euler constant
do not extend unchanged. The positive double kernel, critical
Hurwitz identities, and rational fractional-order certificates give
separate tools on and below the transition. The remaining conjectures
are stated with their proved special cases and their unresolved parts
kept distinct.

\clearpage
\addcontentsline{toc}{section}{References}
\begin{thebibliography}{9}
\bibitem{repo}
ProveIt Contributors,
\emph{Polylogarithms and their Arithmetic Bridges}, canonical manuscript
and accompanying source, repository \texttt{VladimirReshetnikov/ProveIt},
revision \basecommit.
\url{https://github.com/VladimirReshetnikov/ProveIt/tree/28357e8ca63dd78327db91d9be239d75e4462879/Analysis/Polylogarithms/docs/manuscript}.

\bibitem{angular}
ProveIt Contributors,
\emph{Gaussian Polylogarithms Beyond Parity: Exact relation modules,
an octahedral proof of the mixed harmonic sum, and global angular-zero
geometry}, report source at the same pinned revision.
\url{https://github.com/VladimirReshetnikov/ProveIt/blob/28357e8ca63dd78327db91d9be239d75e4462879/Analysis/Polylogarithms/docs/reports/angular-zero-continuation/article.tex}.
Source blob: \nolinkurl{9c2ec9383debad3b1a1248c13f442a0d850487b5}.

\bibitem{dlmf-hurwitz}
NIST Digital Library of Mathematical Functions,
\S25.11, \emph{Hurwitz Zeta Function}, especially 25.11.25 and 25.11.27.
Version 1.2.8, released 15 September 2026; accessed 9 October 2026.
\url{https://dlmf.nist.gov/25.11}.

\bibitem{dlmf-polylog}
NIST Digital Library of Mathematical Functions,
\S25.12, \emph{Polylogarithms}, version 1.2.8.
\url{https://dlmf.nist.gov/25.12}.

\bibitem{dlmf-hyperbolic}
NIST Digital Library of Mathematical Functions,
\S4.36, \emph{Infinite Products and Partial Fractions}, especially 4.36.3,
version 1.2.8.
\url{https://dlmf.nist.gov/4.36}.

\bibitem{liu-pego}
J.-G. Liu and R. L. Pego,
\emph{On generating functions of Hausdorff moment sequences},
Transactions of the American Mathematical Society \textbf{368} (2016),
8499--8518. The arXiv version 4 (12 November 2025) includes a corrigendum.
\url{https://arxiv.org/abs/1401.8052v4}.
\end{thebibliography}
\end{document}
